\scanpage{186}
Cas $X\mathrel{\overset{f}{\text{---}}}S$ propre \inserted{lisse} de dim relative 1. On a un diagramme autodual [dualité = symétrie p.r. au centre)

\diagram{\begin{tikzcd}[ampersand replacement=\&,row sep=small,column sep=small]
\&\&0\&\&\\
\&\& R^1f_*(\underline O_{X/S})\arrow[u]\&\&\\
0\arrow[r]\& R^1f_*(\underline O_X)^{(p)}\arrow[r,"\alpha"]\arrow[ur,"F\ (\mathrm{frob})"]\&\mathcal H^1_{\mathrm{DR}}\arrow[r,"\beta"]\arrow[u]\&f_*(\Omega^1_{X/S})^{(p)}\arrow[r]\&0\\
\&\& f_*(\Omega^1_{X/S})\arrow[u,no head]\arrow[ur,"C\ (\text{op. de Cartier})"']\&\&\\
\&\&0\arrow[u]\&\&
\end{tikzcd}}
\margin{\uline{ligne et colonne \unsure{ex.ct|ex.cte}}}

\margin{$\ast$) $\mathcal H^1_{\mathrm{DR}}=\mathcal H^1_{\mathrm{DR}}(X/S)=R^1f_*(\Omega^\bullet_{X/S})$}

\margin{$\ast\ast$) Colonne est la filtration de $\mathcal H^1$ par $\mathcal H^{1,0}$ et $\mathcal H^{0,1}$}

\margin{$\ast\ast\ast$) $\alpha$ se déduit par passage au quotient de $F_{\mathcal H^1}:\mathcal H^{1(p)}\mathrel{\text{---}}\mathcal H^1$}

\margin{$\ast\ast\ast\ast$) $\beta={}^{t}\alpha$, déduit de $ {}^{t}F_{\mathcal H^1}$ par restriction}

\margin{$\ast\ast\ast\ast\ast$) $F$ et $C$ sont \deleted{duaux} \inserted{transposés} l'un de l'autre, ce sont la puiss. p.ième et Cartier.}

\deleted{On a $\beta\alpha=0$, et on a équivalence entre les conditions}

\deleted{(i) $\alpha$ injectif}

\deleted{(ii) $\beta$ surjectif}

\deleted{(iii) la suite $0\mathrel{\text{---}}R^1f_*(\underline O_X)^{(p)}\xrightarrow{\alpha}\mathcal H^1_{\mathrm{DR}}\xrightarrow{\beta}f_*(\Omega^1_{X/S})^{(p)}\longrightarrow0$ est exacte}

\deleted{(iii bis) $\operatorname{Ker}\beta=\operatorname{Im}\alpha$ \quad i.e. $\operatorname{Ker}\beta=(\operatorname{Ker}\beta)^\perp$ (on a tjs $\supset$) \quad i.e. $\operatorname{Im}\alpha=(\operatorname{Im}\alpha)^\perp$ (on a tjs $\subset$)}

\deleted{(iii ter) $\dim\operatorname{Ker}\beta=\dim\operatorname{Im}\alpha$}

\deleted{(iv) $\operatorname{rang}\alpha=g$ \quad i.e. $\operatorname{rang}\beta=g$ \quad i.e. $\operatorname{rg}F_{\mathcal H^1}=g$}

\margin{\deleted{\uline{NB} on a tjrs $\left\{\begin{aligned}\operatorname{Ker}\alpha&=(\operatorname{Im}\beta)^\perp\\\operatorname{Im}\alpha&=(\operatorname{Ker}\beta)^\perp\end{aligned}\right.$ \deleted{\illegible{2}}}}
\deleted{Pour ceci, il suffit que}
% archive sheet number: 185
